\subsection{The Ramond-sector index behind the determinant}

\begin{definition}[Ramond-sector protected index]
Let $\mathcal H_{\mathrm{RR}}(S)$ be the Hilbert space of a unitary
$\mathcal N=(4,4)$ sigma model on a K3 surface $S$, in the
Ramond--Ramond sector.  Its central charges are
$$c=\overline c=6.$$
Choose a right-moving Ramond supercharge $Q_R$ with
$$Q_R^2=0,\qquad
\{Q_R,Q_R^\dagger\}=2\left(\overline L_0-\frac{\overline c}{24}\right).$$
The K3 elliptic genus is the protected supertrace
$$
Z_{\mathrm{K3}}(\tau,z)=
\operatorname{Tr}_{\mathcal H_{\mathrm{RR}}(S)}
(-1)^{F_L+F_R}q^{L_0-c/24}\overline q^{\,\overline L_0-\overline c/24}
r^{J_0},
$$
where
$$q=e^{2\pi i\tau},\qquad r=e^{2\pi iz}.$$
The $Q_R$-cohomological BPS space in degree $(n,l)$ is
$$
\mathcal H^{\mathrm{BPS}}_{n,l}(S):=
H^\bullet\!\left(
\{\psi\in\mathcal H_{\mathrm{RR}}(S)\mid
(L_0-c/24)\psi=n\psi,\ J_0\psi=l\psi\},Q_R
\right).
$$
\end{definition}

\begin{theorem}[Right-moving localization]
The protected trace localizes to $Q_R$-cohomology:
$$
Z_{\mathrm{K3}}(\tau,z)=
\sum_{n\geq0,\ l\in\Z}
\sdim\mathcal H^{\mathrm{BPS}}_{n,l}(S)\,q^nr^l.
$$
Equivalently,
$$
\sdim\mathcal H^{\mathrm{BPS}}_{n,l}(S)
=[q^nr^l]\,Z_{\mathrm{K3}}(\tau,z).
$$
\end{theorem}

\begin{proof}
If $\overline L_0-\overline c/24>0$, the identity
$$
\{Q_R,Q_R^\dagger\}=2\left(\overline L_0-\frac{\overline c}{24}\right)
$$
pairs the corresponding states by opposite $(-1)^{F_R}$ parity.  Their
contribution to the supertrace is zero.  The unpaired states are exactly
the harmonic representatives for $Q_R$-cohomology, hence have
$$\overline L_0=\overline c/24.$$
The supertrace is therefore independent of $\overline q$ and equals the
graded superdimension of the displayed cohomology.
\end{proof}

\begin{definition}[Half-K3 index object]
Use the Eichler--Zagier normalization already fixed in the manuscript,
\cite{EZ:1},
$$
Z_{\mathrm{K3}}(\tau,z)=2\phi_{0,1}(\tau,z),\qquad
\phi_{0,1}(\tau,z)=\sum_{n\geq0,\ l\in\Z}f(n,l)q^nr^l.
$$
Thus
$$
\sdim\mathcal H^{\mathrm{BPS}}_{n,l}(S)=2f(n,l).
$$
The class
$$
\mathbb U_{n,l}\in K_0(\mathrm{sVect}_{\mathrm{fd}})
$$
defined by
$$
\sdim\mathbb U_{n,l}=f(n,l)
$$
is the normalized half-index object.  It is not a canonical half of
$\mathcal H^{\mathrm{BPS}}_{n,l}(S)$ as a Hilbert space.  It is the
integral virtual super vector space whose signed dimension is the
Borcherds input coefficient.
\end{definition}

\begin{remark}[Status of the half-index]
The Ramond-sector space
$\mathcal H^{\mathrm{BPS}}_{n,l}(S)$ is a physical protected Hilbert
object.  The class $\mathbb U_{n,l}$ is an index shadow.  Its existence
uses only the integrality of $f(n,l)$:
$$
f(n,l)\geq0:\quad \mathbb U_{n,l}\simeq\C_{\overline0}^{\,f(n,l)},
\qquad
f(n,l)<0:\quad \mathbb U_{n,l}\simeq\C_{\overline1}^{\,-f(n,l)}
$$
is one possible representative.  No construction of a canonical
square-root Hilbert space is asserted here.
\end{remark}

\begin{theorem}[DMVV index and normalized shadow]
Let
$$
Z_{\mathrm{K3}}(\tau,z)=\sum_{n\geq0,\ l\in\Z}c(n,l)q^nr^l,
\qquad c(n,l)=2f(n,l).
$$
Dijkgraaf--Moore--Verlinde--Verlinde \cite{DMVV} prove the
second-quantized elliptic-genus identity
$$
\sum_{N\geq0}p^N Z_{\mathrm{ell}}\!\left(\operatorname{Sym}^N S;\tau,z\right)
=
\prod_{\substack{m>0,\ n\geq0\\ l\in\Z}}
\left(1-p^m q^n r^l\right)^{-c(nm,l)}.
$$
Replacing $c(n,l)$ by $f(n,l)$ gives the normalized virtual coefficient
system used by the Borcherds lift:
$$
\operatorname{PE}\!\left(\sum_{\substack{m>0,\ n\geq0\\ l\in\Z}}
f(nm,l)p^mq^nr^l\right)
=
\prod_{\substack{m>0,\ n\geq0\\ l\in\Z}}
\left(1-p^m q^n r^l\right)^{-f(nm,l)}.
$$
\end{theorem}

\begin{remark}[Boundary factors]
The DMVV theorem accounts for the sector $m>0$.  The Borcherds product
for $\Delta_5$ uses the larger positive semigroup
$$
\Gamma_{\mathrm{eff}}=
\{(n,l,m)\mid m>0,\ n\geq0,\ l\in\Z\}
\cup
\{(n,l,0)\mid n>0,\ l\in\Z\}
\cup
\{(0,l,0)\mid l<0\}.
$$
The $m=0$ terms are cusp and Weyl-chamber boundary terms of the
automorphic product, not additional Hilbert spaces supplied by DMVV.
They may be encoded in the same virtual index notation because the
determinant only reads the numbers $f(0,l)$.
\end{remark}

\begin{definition}[One-particle virtual super vector space]
For $\gamma=(n,l,m)\in\Gamma_{\mathrm{eff}}$ put
$$x^\gamma=q^nr^ls^m.$$
Define a $\Gamma_{\mathrm{eff}}$-graded virtual super vector space
$$
\mathcal V=\bigoplus_{\gamma\in\Gamma_{\mathrm{eff}}}\mathcal V_\gamma
$$
in $K_0(\mathrm{sVect}_{\mathrm{fd}})$ by
$$
\mathcal V_{n,l,m}:=\mathbb U_{nm,l},\qquad
\sdim\mathcal V_{n,l,m}=f(nm,l).
$$
This is the precise sense in which the manuscript uses a one-particle
BPS space.  It is a protected-index object.  It is not, by itself, a
construction of all microscopic BPS states, nor a construction of the
Lie bracket of $\mathfrak g_{\Delta_5}$.
\end{definition}

\begin{lemma}[Determinants see only signed dimensions]
Let
$$V=V_{\overline0}\oplus V_{\overline1}$$
be a finite-dimensional super vector space, and let $x$ be a scalar
formal monomial.  Then
$$
\sdet_V(1-x)
=
\frac{\det(1-x\mid V_{\overline0})}
{\det(1-x\mid V_{\overline1})}
=
(1-x)^{\dim V_{\overline0}-\dim V_{\overline1}}
=(1-x)^{\sdim V}.
$$
Consequently two virtual super vector spaces with the same signed
dimension define the same determinant factor.
\end{lemma}

\begin{proof}
The endomorphism $1-x$ acts by the same scalar on every homogeneous
basis vector.  The even basis vectors contribute $\det(1-x)$ to the
numerator, the odd basis vectors contribute $\det(1-x)$ to the
denominator.  The exponent is therefore
$$\dim V_{\overline0}-\dim V_{\overline1}=\sdim V.$$
\end{proof}

\begin{theorem}[Index determinant]
For the virtual one-particle object above, the formal Fock character is
$$
\sch\mathcal F(\mathcal V)
=
\prod_{\gamma\in\Gamma_{\mathrm{eff}}}
\left(1-x^\gamma\right)^{-\sdim\mathcal V_\gamma}
=
\prod_{(n,l,m)\in\Gamma_{\mathrm{eff}}}
\left(1-q^nr^ls^m\right)^{-f(nm,l)},
$$
and
$$
\sdet_{\mathcal V}(1-x)
=
\prod_{\gamma\in\Gamma_{\mathrm{eff}}}
\left(1-x^\gamma\right)^{\sdim\mathcal V_\gamma}
=
\prod_{(n,l,m)\in\Gamma_{\mathrm{eff}}}
\left(1-q^nr^ls^m\right)^{f(nm,l)}.
$$
Hence the determinant used in the manuscript is
$$
\mathcal D_X(Z)
=64q^{1/2}r^{1/2}s^{1/2}\sdet_{\mathcal V}(1-x)
$$
$$
=64q^{1/2}r^{1/2}s^{1/2}
\prod_{(n,l,m)\in\Gamma_{\mathrm{eff}}}
\left(1-q^nr^ls^m\right)^{f(nm,l)}.
$$
\end{theorem}

\begin{proof}
For one even generator of charge $\gamma$,
$$\sch\operatorname{Sym}(\C_{\overline0}x^\gamma)
=(1-x^\gamma)^{-1}.$$
For one odd generator of charge $\gamma$,
$$\sch\bigwedge(\C_{\overline1}x^\gamma)=1-x^\gamma.$$
Multiplication over a homogeneous basis gives the Fock character, and
the preceding lemma gives the inverse superdeterminant.  Since each
factor depends only on $\sdim\mathcal V_\gamma=f(nm,l)$, the determinant
is well-defined for the virtual class $\mathcal V_\gamma$.
\end{proof}

\begin{assumption}[Genuine chiral square-root Hilbert lift]
A stronger physical assertion is the existence of an honest
one-particle Hilbert space
$$
\mathcal H^{\mathrm{1p,ch}}=
\bigoplus_{\gamma\in\Gamma_{\mathrm{eff}}}
\mathcal H^{\mathrm{1p,ch}}_\gamma
$$
whose protected superdimensions are
$$
\operatorname{Tr}_{\mathcal H^{\mathrm{1p,ch}}_\gamma}(-1)^F
=f(nm,l),\qquad \gamma=(n,l,m).
$$
This is an additional physical lift.  DMVV proves the
second-quantized index for the full K3 elliptic genus.  It does not
construct a canonical Hilbert-space square root realizing
$\phi_{0,1}=Z_{\mathrm{K3}}/2$ state by state.
\end{assumption}

\begin{remark}[Relation to BPS algebras]
Harvey--Moore \cite{HM} give the physical framework in which BPS
states carry charge gradings and form algebraic structures.  The
Dijkgraaf--Verlinde--Verlinde dyon-counting formula \cite{DVV} explains
where Igusa-type products occur in the $N=4$ BPS index problem.
The present determinant statement uses only the protected supertrace and
the Borcherds product coefficients.  It proves equality of signed
graded dimensions.  It does not prove an isomorphism between a
microscopic BPS Hilbert algebra and $\mathfrak g_{\Delta_5}$.
\end{remark}
